% GENERATED FILE. Edit canonical lessons, then run npm run generate:books.
% canonical-source-hash: fnv1a64:294475b69d2a69f5
% canonical-lessons: ZH-C246-cunzi, ZH-C246-qiao, ZH-C246-shitou, ZH-C246-chitang, ZH-C246-huachuan

\chapter{Out in the Hills}
\label{ch:zh-out-in-the-hills}
\hlchaptermodality{246}

\begin{chapteropening}
\textbf{By the end of this chapter:} \emph{I can talk about a village, a bridge, stones, a pond and rowing a boat.}

Complete the last lesson of chapter 246: I can talk about a village, a bridge, stones, a pond and rowing a boat.
\end{chapteropening}

\section[cūnzi]{cūnzi — village}

\label{lesson:ZH-C246-cunzi}



\begin{quote}
Before the new one: say the Mandarin for mood, then the Mandarin for to look after yourself.
\end{quote}

\subsection*{You'll want to know: cūnzi}
\emph{cūnzi} — \textquotedblleft{}village\textquotedblright{}.

\begin{quote}
\emph{zhège cūnzi hěn xiǎo} — This village is tiny.
\end{quote}

Its characters come later; for now it is a word for the ear.

\begin{cousinweb}
\emph{cūnzi} is built from \emph{cūn} 'village' and the noun ending \zh{子}, both inherited from Middle Chinese, \emph{tshwon} and \emph{tsiX} (Guangyun, in Baxter's transcription).
\end{cousinweb}

\begin{grammarlens}[title={What you've built}]
One more word for saying what I do.
\end{grammarlens}

\subsection*{Your turn}
\begin{itemize}
\raggedright
  \item \emph{Say it:} \emph{cūnzi}
  \item \emph{Say it:} \emph{cūnzi}, once more
  \item \emph{Say it:} \emph{zhège cūnzi hěn xiǎo} — This village is tiny.
  \item \emph{Recall:} say the Mandarin for mood, then the Mandarin for to look after yourself, then say \emph{cūnzi} again
\end{itemize}

\begin{tcolorbox}[breakable,colback=teal!4,colframe=teal!35!black,title={Before you move on}]
What does cūnzi mean? (\textquotedblleft{}Village\textquotedblright{}.) Say it once more.
\end{tcolorbox}
\section[qiáo]{qiáo — bridge}

\label{lesson:ZH-C246-qiao}



\begin{quote}
Before the new one: say the Mandarin for to look after yourself, then the Mandarin for a village.
\end{quote}

\subsection*{You'll want to know: qiáo}
\emph{qiáo} — \textquotedblleft{}bridge\textquotedblright{}.

\begin{quote}
\emph{cūnzi lǐ de qiáo hěn lǎo} — The bridge in the village is very old.
\end{quote}

Its characters come later; for now it is a word for the ear.

\begin{cousinweb}
\emph{qiáo} is inherited from Middle Chinese \emph{gjew} (Guangyun, in Baxter's transcription), glossed 'a beam over water'.
\end{cousinweb}

\begin{grammarlens}[title={What you've built}]
One more word for saying what I do.
\end{grammarlens}

\subsection*{Your turn}
\begin{itemize}
\raggedright
  \item \emph{Say it:} \emph{qiáo}
  \item \emph{Say it:} \emph{qiáo}, once more
  \item \emph{Say it:} \emph{cūnzi lǐ de qiáo hěn lǎo} — The bridge in the village is very old.
  \item \emph{Recall:} say the Mandarin for to look after yourself, then the Mandarin for a village, then say \emph{qiáo} again
\end{itemize}

\begin{tcolorbox}[breakable,colback=teal!4,colframe=teal!35!black,title={Before you move on}]
What does qiáo mean? (\textquotedblleft{}Bridge\textquotedblright{}.) Say it once more.
\end{tcolorbox}
\section[shítou]{shítou — stone; rock}

\label{lesson:ZH-C246-shitou}



\begin{quote}
Before the new one: say the Mandarin for a village, then the Mandarin for a bridge.
\end{quote}

\subsection*{You'll want to know: shítou}
\emph{shítou} — \textquotedblleft{}stone; rock\textquotedblright{}.

\begin{quote}
\emph{nà kuài shítou hěn dà} — That rock is huge.
\end{quote}

Its characters come later; for now it is a word for the ear.

\begin{cousinweb}
\emph{shítou} is built from \emph{shí} 'stone' and \zh{头}, here a light noun ending, both inherited from Middle Chinese, \emph{dzyek} and \emph{duw} (Guangyun, in Baxter's transcription), which explains the first, after an older book, as the body of a mountain.
\end{cousinweb}

\begin{grammarlens}[title={What you've built}]
One more word for saying what I do.
\end{grammarlens}

\subsection*{Your turn}
\begin{itemize}
\raggedright
  \item \emph{Say it:} \emph{shítou}
  \item \emph{Say it:} \emph{shítou}, once more
  \item \emph{Say it:} \emph{nà kuài shítou hěn dà} — That rock is huge.
  \item \emph{Recall:} say the Mandarin for a village, then the Mandarin for a bridge, then say \emph{shítou} again
\end{itemize}

\begin{tcolorbox}[breakable,colback=teal!4,colframe=teal!35!black,title={Before you move on}]
What does shítou mean? (\textquotedblleft{}Stone; rock\textquotedblright{}.) Say it once more.
\end{tcolorbox}
\section[chítáng]{chítáng — pond}

\label{lesson:ZH-C246-chitang}



\begin{quote}
Before the new one: say the Mandarin for a bridge, then the Mandarin for a stone.
\end{quote}

\subsection*{You'll want to know: chítáng}
\emph{chítáng} — \textquotedblleft{}pond\textquotedblright{}.

\begin{quote}
\emph{cūnzi lǐ yǒu yí ge chítáng} — There's a pond in the village.
\end{quote}

Its characters come later; for now it is a word for the ear.

\begin{cousinweb}
\emph{chítáng} is built from \emph{chí} 'a pool' and \emph{táng} 'a bank'. The dictionary (Guangyun, in Baxter's transcription) explains the first, \emph{drje} there, as standing water, and glosses the second, \emph{dang} there, with a word for a dyked pond.
\end{cousinweb}

\begin{grammarlens}[title={What you've built}]
One more word for saying what I do.
\end{grammarlens}

\subsection*{Your turn}
\begin{itemize}
\raggedright
  \item \emph{Say it:} \emph{chítáng}
  \item \emph{Say it:} \emph{chítáng}, once more
  \item \emph{Say it:} \emph{cūnzi lǐ yǒu yí ge chítáng} — There's a pond in the village.
  \item \emph{Recall:} say the Mandarin for a bridge, then the Mandarin for a stone, then say \emph{chítáng} again
\end{itemize}

\begin{tcolorbox}[breakable,colback=teal!4,colframe=teal!35!black,title={Before you move on}]
What does chítáng mean? (\textquotedblleft{}Pond\textquotedblright{}.) Say it once more.
\end{tcolorbox}
\section[huáchuán]{huáchuán — to row a boat}

\label{lesson:ZH-C246-huachuan}



\begin{quote}
Before the new one: say the Mandarin for a stone, then the Mandarin for a pond.
\end{quote}

\subsection*{You'll want to know: huáchuán}
\emph{huáchuán} — \textquotedblleft{}to row a boat\textquotedblright{}.

\begin{quote}
\emph{zhōumò wǒmen qù gōngyuán huáchuán} — At weekends we go rowing in the park.
\end{quote}

Its characters come later; for now it is a word for the ear.

\begin{cousinweb}
\emph{huáchuán} is built from \emph{huá} 'to pole along' and \emph{chuán} 'boat', both inherited from Middle Chinese, \emph{hwae} and \emph{zywen} (Guangyun, in Baxter's transcription), which glosses the first 'to push a boat forward'.
\end{cousinweb}

\begin{grammarlens}[title={What you've built}]
One more word for saying what I do.
\end{grammarlens}

\subsection*{Your turn}
\begin{itemize}
\raggedright
  \item \emph{Say it:} \emph{huáchuán}
  \item \emph{Say it:} \emph{huáchuán}, once more
  \item \emph{Say it:} \emph{zhōumò wǒmen qù gōngyuán huáchuán} — At weekends we go rowing in the park.
  \item \emph{Recall:} say the Mandarin for a stone, then the Mandarin for a pond, then say \emph{huáchuán} again
\end{itemize}

\begin{tcolorbox}[breakable,colback=teal!4,colframe=teal!35!black,title={Before you move on}]
What does huáchuán mean? (\textquotedblleft{}To row a boat\textquotedblright{}.) Say it once more.
\end{tcolorbox}
